% ─────────────────────────────────────────────────────────────────────────────
\subsection*{Overview}

This section derives an IA-DEVS formulation of the CKW10 Load Balancer
Model, starting from the STDEVS specification in Section~\ref{sec:ckw10-lbm}.
The derivation proceeds in two steps.
\begin{enumerate}
  \item \textbf{STDEVS $\to$ UA-DEVS.}
        Each stochastic transition function $\delta(s,r)$ with $r \sim U(0,1)$
        is replaced by the set of all values it can produce:
        $\{\delta(s,r) \mid r \in (0,1)\}$.  This set is an uncertain set
        in the sense of UA-DEVS~\cite{VDW21}.
  \item \textbf{UA-DEVS $\to$ IA-DEVS.}
        IA-DEVS approximates each UA-DEVS uncertain \emph{set} with an
        \emph{interval} and computes exclusively with its two endpoint
        values.  This is what makes IA-DEVS computable: any simulation step
        reduces to a finite number of endpoint evaluations regardless of
        how many values lie inside the interval.
        Note that the interval need not be bounded: $(0,+\infty)$ is as
        computable as $[0.01, 0.60]$ because only the endpoints are ever
        evaluated.

        The set-to-interval mapping for each LBM component is:
        \begin{itemize}
          \item LG and servers: the UA-DEVS set $(0,+\infty)$ is
            represented directly as the IA-DEVS interval $(0,+\infty)$,
            the exact support of the exponential distribution.
            The resulting loss bound is $P_{\mathrm{loss}} \in [0,1]$,
            which is trivially correct; a tighter interval such as a
            quantile-truncated $[\sigma_L, \sigma_U]$ yields more
            informative bounds at the cost of excluding tail trajectories.
          \item WB: the UA-DEVS set $\{1,2\}$ is approximated by the
            IA-DEVS interval $[1,2]$.  In general IA-DEVS intervals
            over-approximate the underlying set; here the approximation
            is exact for integer outputs since no integers lie strictly
            between $1$ and $2$.
        \end{itemize}
\end{enumerate}

% ─────────────────────────────────────────────────────────────────────────────
\subsection*{Step 1: STDEVS to UA-DEVS Translation}

The translation looks only at the \emph{observable outputs and reachable
states} of each model component, not at the internal random mechanism used
to select among them.  The specific distribution (exponential, uniform,
Bernoulli, \ldots) is an implementation detail of the STDEVS model and
does not appear in the UA-DEVS specification.

\subsubsection*{Load Generator (LG)}

The LG specification says: after each internal transition, the next
inter-departure time is some positive value.  The set of all reachable
next states is therefore
\[
  S^{LG}_{\text{next}} = (0,\, +\infty) \subset \mathbb{R}^+.
\]
The output is deterministic: always $(\mathit{task}_1,\, \texttt{out}_1)$.

\subsubsection*{Weighted Balancer (WB)}

The WB specification says: each arriving task is routed to
\emph{either} $\texttt{out}_1$ (server~$S1$) or $\texttt{out}_2$
(server~$S2$) — both are possible outcomes.
Encoding the output port as an integer
($\texttt{out}_1 \mapsto 1,\; \texttt{out}_2 \mapsto 2$),
the UA-DEVS uncertain output is the set $\{1,\, 2\} \subset \mathbb{Z}$.
The time advance after routing is $+\infty$ (WB becomes passive),
which is deterministic.

\subsubsection*{Servers S1 and S2}

The server specification says: when the server is \emph{idle} and a
task arrives, it begins service for some positive duration.  The set
of all reachable service-time states is
\[
  S^{Si}_{\text{next}} = (0,\, +\infty) \subset \mathbb{R}^+.
\]
When the server is \emph{busy} and a task arrives, the task is dropped;
this branch is deterministic and is preserved unchanged in UA-DEVS.

% ─────────────────────────────────────────────────────────────────────────────
\subsection*{Step 2: UA-DEVS to IA-DEVS Bounding}

\paragraph{WB routing over-approximation.}
For tractability we decompose the two-server LBM into two independent
$M/M/1/1$ sub-models, one per server, each with arrival rate
$\lambda_s = b_f d_r$.  This is exact when $b_f = 0.5$ and $\mu_1 = \mu_2$
because the routing is symmetric.  In the IA-DEVS formulation, routing
uncertainty is handled by applying the same interval bounds to both
sub-models independently.

% ─────────────────────────────────────────────────────────────────────────────
\subsection*{IA-DEVS Single-Server Model ($M/M/1/1$)}

We formulate the per-server sub-model as two IA-DEVS atomic components
connected in series: an \textit{ia\_load\_generator} (LG-IA) and an
\textit{ia\_mm11\_server} (S-IA).

\subsubsection*{LG-IA: Load Generator (IA-DEVS)}

State $\tilde\sigma \in \mathcal{I}(\mathbb{R}^+_0)$ (elapsed time since
last departure; passive state = empty interval).

\begin{align*}
  \mathrm{TA}^{\mathrm{LG}}(\tilde\sigma)
    &= (0,\,+\infty) \\
  \Delta^{\mathrm{LG}}_{\mathrm{int}}(\tilde\sigma)
    &= [0,\,0] \\
  \Delta^{\mathrm{LG}}_{\mathrm{ext}}(\tilde\sigma,\,\tilde e,\,\tilde x)
    &= \tilde\sigma + \tilde e
       \quad\text{(accumulate elapsed; no meaningful input)} \\
  \Lambda^{\mathrm{LG}}(\tilde\sigma)
    &= [1,\,1]
\end{align*}

The time-advance interval $(0,+\infty)$ is the exact support of
$\mathrm{Exp}(\lambda_s)$.  Because the exponential distribution is
memoryless, the remaining time until the next departure is always
$(0,+\infty)$ regardless of elapsed time, so $\mathrm{TA}^{\mathrm{LG}}$
does not depend on $\tilde\sigma$.

\subsubsection*{S-IA: $M/M/1/1$ Server (IA-DEVS)}

State: $\emptyset$ (idle/passive) or $\tilde s \in \mathcal{I}(\mathbb{R}^+_0)$
(remaining service time when busy).

\begin{align*}
  \mathrm{TA}^S(\emptyset)
    &= \emptyset \quad\text{(passive)} \\
  \mathrm{TA}^S(\tilde s)
    &= \tilde s \\
  \Delta^S_{\mathrm{int}}(\tilde s)
    &= \emptyset \quad\text{(service completes; go idle)} \\
  \Delta^S_{\mathrm{ext}}(\emptyset,\,\tilde e,\,\tilde x)
    &= (0,\,+\infty)
       \quad\text{(accept task; start service)} \\
  \Delta^S_{\mathrm{ext}}(\tilde s,\,\tilde e,\,\tilde x)
    &= \max(0,\; \tilde s - \tilde e)
       \quad\text{(drop task; update remaining time)} \\
  \Lambda^S(\tilde s)
    &= [1,\,1]
       \quad\text{(completed task signal)}
\end{align*}

A \emph{loss event} is observed when $\Delta^S_{\mathrm{ext}}$ is called
on a non-empty state, indicating that the server was busy at arrival time.

% ─────────────────────────────────────────────────────────────────────────────
\subsection*{Analytical $P_{\mathrm{loss}}$ Bounds}

The Erlang~B loss probability $P_{\mathrm{loss}}(\rho) = \rho/(1+\rho)$
is monotone increasing in the traffic intensity $\rho = \lambda_s \cdot s$.

With service time $s \in (0,+\infty)$ and fixed arrival rate $\lambda_s > 0$,
the traffic intensity satisfies $\rho \in (0,+\infty)$.  By monotonicity of
$P_{\mathrm{loss}}(\rho)$ and taking the closure of the limit values
$\rho \to 0^+$ and $\rho \to +\infty$:
\[
  \tilde P_{\mathrm{loss}} \;\in\; [0,\, 1].
\]

The exact STDEVS value from the Erlang~B formula is
$P_{\mathrm{loss}}^* = \rho^*/(1+\rho^*) = 1/2 = 0.5$
(since $\rho^* = \lambda_s/\mu = 5/5 = 1$),
which satisfies $0 \leq 0.5 \leq 1$ as required. \hfill$\square$

\paragraph{Remark.}
The $[0,1]$ bound is trivially correct but uninformative.  A tighter bound
can be obtained by truncating the exponential support to a quantile interval
$[q_{p_L},\, q_{1-p_L}]$ for chosen tail probability $p_L$, where
$q_\alpha(\lambda) = -\tfrac{1}{\lambda}\ln(1-\alpha)$.
For example, with $p_L = 0.05$ (90\,\% central region) and
$\lambda_s = \mu = 5$, the bounds
$[\sigma_L,\sigma_U] = [s_L,s_U] \approx [0.0103,\; 0.5992]\,\text{s}$
yield $\tilde P_{\mathrm{loss}} \approx [0.017,\; 0.983]$.
As $p_L \to 0.5$ (point interval at the median), the bound collapses to
the exact value but the over-approximation guarantee is lost.

% ─────────────────────────────────────────────────────────────────────────────
\subsection*{Feasibility Note on BFS Simulation}

A full IA-DEVS BFS simulation of the LBM at operational scale
($d_r = 10$, $T = 10\,000$\,s, ${\sim}10^5$ events) is not computationally
feasible: whenever the LG time-advance interval $(0,+\infty)$ and the
server remaining-service interval overlap, each arrival creates an ordering
branch (``arrival before completion'' vs.\ ``arrival after completion''),
leading to exponential branching in the number of events.  With
$(0,+\infty)$ intervals the two ranges always overlap, so every event
produces a branch.

The analytical bound above is the IA-DEVS result without BFS, derived by
applying the bounding functions directly to the Erlang~B formula.  The
C\texttt{++} implementation below verifies this bound and demonstrates that
the IA-DEVS atomic models are correctly specified, including a short
5-event simulation trace that shows the interval widening effect.

% ─────────────────────────────────────────────────────────────────────────────
\subsection*{5-Event Simulation Trace}

Initial state: LG-IA elapsed $\tilde\sigma_0 = [0,0]$; S-IA state $= \emptyset$.

\begin{center}
\footnotesize
\begin{tabular}{@{}clll@{}}
\toprule
Event & Action & S-IA remaining $\tilde s$ & Loss? \\
\midrule
1 & LG fires; S-IA idle $\to$ accept & $(0,\,+\infty)$ & no \\
2 & LG fires (elapsed $\tilde e \in (0,+\infty)$);
    S-IA: $\max(0,\,\tilde s - \tilde e)$
  & $[0,\,+\infty)$ & possible \\
3 & LG fires; S-IA: $\max(0,\,\tilde s - \tilde e)$
  & $[0,\,+\infty)$ & possible \\
4 & LG fires & $[0,\,+\infty)$ & possible \\
5 & LG fires & $[0,\,+\infty)$ & possible \\
\bottomrule
\end{tabular}
\end{center}

After the first acceptance the remaining service interval is $(0,+\infty)$.
At event~2, interval subtraction gives
$(0,+\infty) - (0,+\infty) = (-\infty,+\infty)$; clamping the lower bound
to~$0$ (closed) yields $[0,+\infty)$.  This state is absorbing: subsequent
elapsed-time subtractions leave $[0,+\infty)$ unchanged.  The closed
lower bound $0$ means the server \emph{could} have already finished
(possible non-loss); the unbounded upper means it could still be running
(possible loss).  The exact STDEVS simulation determines which branch
actually occurs.
